\documentclass[10pt,a4paper]{amsart}

\usepackage[T1]{fontenc}
\usepackage{amsmath,amssymb,amsthm,mathtools}
\usepackage[a4paper,left=18mm,right=18mm,top=15mm,bottom=15mm]{geometry}
\usepackage[hidelinks]{hyperref}
\usepackage{microtype}
\microtypesetup{expansion=false}

\newtheorem{theorem}{Theorem}[section]
\newtheorem{lemma}[theorem]{Lemma}
\newtheorem{corollary}[theorem]{Corollary}
\theoremstyle{remark}
\newtheorem{remark}[theorem]{Remark}

\newcommand{\Z}{\mathbb Z}
\newcommand{\Q}{\mathbb Q}
\newcommand{\R}{\mathbb Z[1/2]}
\newcommand{\E}{\mathbb Q/\mathbb Z[1/2]}
\newcommand{\U}{\operatorname U}
\newcommand{\Sym}{\operatorname {Sym}}
\newcommand{\aug}{\varpi}
\newcommand{\dlt}{\delta}
\newcommand{\gm}{\gamma}

\title[Odd torsion in the fifth Lie dimension quotient]
{Odd torsion in the fifth Lie dimension quotient:\\
a four-weight proof}
\author{}
\date{22 August 2026}

\begin{document}

\begin{abstract}
For a Lie ring $L$, let
$\delta_n(L)=L\cap\varpi(L)^n$, where $\varpi(L)$ is the augmentation
ideal of its enveloping algebra.  We prove that
$\delta_5(L)/\gamma_5(L)$ has no odd torsion.  Localization at $2$
reduces the theorem directly to $\delta_5(K)=0$ for a class-four
$\mathbb Z[1/2]$-Lie algebra $K$.  Its metabelian quotient puts a
possible element in $K''$.  There are then only two exposed PBW edges
in bracket weight four, of shapes $211$ and $31$.  Their readouts are
explicit Koszul coboundaries.  Polarizing only the two weight-one
entries makes their common $22$ root the weight-four face of one full
commutator relation, so it cancels without a separate overlap ledger.
Only division by $2$ is used.
\end{abstract}

\maketitle

For a Lie ring $A$ write $\U(A)$ for its universal enveloping algebra,
$\aug(A)$ for its augmentation ideal, and
\[
 \dlt_s(A)=A\cap\aug(A)^s,
 \qquad \gm_{s+1}(A)=[\gm_s(A),A].
\]
We use the established metabelian estimate
\begin{equation}
                 2\dlt_s(M)\subseteq\gm_s(M)
                 \qquad(M''=0),                         \label{eq:PS}
\end{equation}
proved in \cite[Theorem 2.1]{PS}.  Intersections are inverse images
under the canonical map $A\to\U(A)$.

\section{The class-four packet}

Throughout this section scalars are in $\R$, and a character means an
$\R$-linear map to $\E$.  Multiplication by $2$ is invertible on $\E$;
moreover $\E$
detects nonzero elements of finitely generated $\R$-modules (use the
structure theorem over the PID $\R$).  We prove the only new ingredient
first.  Dimension subrings here are formed in $\U_{\R}$.  This causes no
conflict with \eqref{eq:PS}.  Indeed, for $\R$-modules $P,Q$ the map
$P\otimes_{\Z}Q\to P\otimes_{\R}Q$ is an isomorphism: a factor $1/2$
may be moved from one entry to the other.  Hence the nonunital tensor
algebras over $\Z$ and $\R$ agree.  Quotienting by the same enveloping
relations therefore identifies the augmentation ideals of
$\U_{\Z}(K)$ and $\U_{\R}(K)$, with all their powers.  Their lower
central terms agree as well, since $r[x,y]=[rx,y]$ for $r\in\R$.

\begin{lemma}[four-weight packet]\label{lem:packet}
Let $K$ be an $\R$-Lie algebra with $\gm_5(K)=0$.  Then
\[
                         \dlt_5(K)\cap K''=0.
\]
\end{lemma}

\begin{proof}
It is enough to consider finitely generated $K$.  Indeed, a witness
for $a\in\delta_5(K)$ is a finite tensor-algebra identity.  Adjoin to
its entries finitely many entries of an expression for $a\in K''$.
The $\R$-subalgebra they generate is still of class at most four, is
finitely generated as an $\R$-module (it is a quotient of a
finite-rank free nilpotent algebra of class four), and contains the
same two witnesses.

Suppose that $0\ne a\in\delta_5(K)\cap K''$.  Put
\[
 V=K/\gm_2K,\quad B=\gm_2K/\gm_3K,\quad
 W=\gm_3K/\gm_4K,\quad D=\gm_4K.
\]
Since $K''=[\gamma_2K,\gamma_2K]\subseteq D$, the finitely generated
module $D$ has a character
$\chi:D\to\E$ with $\chi(a)\ne0$.  Bracketing gives a well-defined map
\begin{equation}
                 g:V\otimes W\longrightarrow D,
                 \qquad g(x,w)=[x,w].                  \label{eq:g}
\end{equation}
Changing its lifts changes the bracket by
$[\gamma_2K,\gamma_3K]+[K,\gamma_4K]\subseteq\gamma_5K=0$.
For $b\in B$ and $x,y\in V$ define
\begin{equation}
 \Theta^+_\chi(b;xy)=\frac12\chi
 \bigl([[b,x],y]+[[b,y],x]\bigr).                     \label{eq:theta}
\end{equation}
This is independent of all lifts: changing a $B$-lift by $\gamma_3K$
or a $V$-lift by $\gamma_2K$ changes the bracket by $\gamma_5K$.
It is symmetric in $x,y$, and Jacobi gives
\begin{equation}
 [[b,x],y]-[[b,y],x]=[b,[x,y]].                         \label{eq:jacobi}
\end{equation}

Choose a finite presentation $K=F/\mathcal R$, where
$F=F_1\oplus\cdots\oplus F_4$ is the free nilpotent $\R$-Lie algebra
of class four on finitely many generators.  The following additive
presentations have free terms, since $\R$ is a PID:
\[
 \begin{array}{lll}
 S_1=\ker(F_1\oplus F_2\to K/\gamma_3K),
   &S_0=F_1\oplus F_2,&\partial_S:S_1\hookrightarrow S_0,\\
 R_1=\ker(F_1\oplus F_2\oplus F_3\to K/\gamma_4K),
   &R_0=F_1\oplus F_2\oplus F_3,&\partial_R:R_1\hookrightarrow R_0.
 \end{array}
\]
For $s\in S_0$ write $s_V,s_B$ for the images of its $F_1,F_2$
coordinates in $V,B$; for $z\in R_0$ define $z_V,z_W$ similarly.
On symmetric powers set
\begin{align}
 T_3(s_1s_2s_3)
   &=\Theta^+_\chi((s_1)_B;(s_2)_V(s_3)_V)
     +\Theta^+_\chi((s_2)_B;(s_1)_V(s_3)_V)\notag\\
   &\qquad+\Theta^+_\chi((s_3)_B;(s_1)_V(s_2)_V),       \label{eq:T3}\\
 T_2(z_1z_2)
   &=-\chi g((z_1)_V,(z_2)_W)
     -\chi g((z_2)_V,(z_1)_W).                         \label{eq:T2}
\end{align}
These formulas are symmetric, hence well defined.  If
$K_m(C)_1=C_1\otimes\Sym^{m-1}C_0$ and
$\partial(u;v)=\partial(u)v$, put $dT=T\circ\partial$.  Then
\begin{align}
 (dT_3)(U;vw)&=\Theta^+_\chi(U_B;v_Vw_V),               \label{eq:dT3}\\
 (dT_2)(Y;z)&=-\chi g(z_V,Y_W).                         \label{eq:dT2}
\end{align}
Indeed $U_V=Y_V=0$, so the displayed terms are the only survivors.
Also $T_3$ and $T_2$ are supported respectively in weights
$2+1+1=4$ and $3+1=4$.

We now extract the two chains to which
\eqref{eq:dT3}--\eqref{eq:dT2} apply.  Choose
$f\in F''\cap F_4$ mapping to $a$; this is possible because
$F''\subseteq F_4$ in class four and $F''\twoheadrightarrow K''$.
Since $a\in\delta_5(K)$, there is
$\theta\in\mathcal R\U(F)$ such that
\begin{equation}
                         \theta-f\in\aug(F)^5.          \label{eq:witness}
\end{equation}
Here the two-sided relation ideal is $\mathcal R\U(F)$, since
$x\rho=\rho x+[x,\rho]$ and $[x,\rho]\in\mathcal R$.
Write $\theta$ as a finite sum of rows $\rho u$, giving every occurrence
of the full relation $\rho$ its own label.  Fix homogeneous PBW
bases and work modulo total weight at least five.  Put
$\rho^{[k]}=\rho_1+\cdots+\rho_k$ for the homogeneous pieces of
$\rho$, and write $\langle\rho,k\rangle$ for an occurrence which
evaluates as $\rho^{[k]}$ but retains the label $\rho$.  Besides the
ordinary PBW swap, use the literal labelled rules
\begin{align}
 vu&=uv+[v,u],\notag\\
 v\langle\rho,k\rangle
   &=\langle\rho,k\rangle v+
     \langle[v,\rho],k+h\rangle &&(v\in F_h),\notag\\
 \langle\rho,k\rangle
   &=\langle\rho,k-1\rangle+\rho_k .                  \label{eq:collect}
\end{align}
Indices are capped at four.  The second equality is exact in the
range retained because
$\operatorname {pr}_{\leq k+h}[v,\rho]=[v,\rho^{[k]}]$.
Thus its correction is still labelled by the full relation
$[v,\rho]\in\mathcal R$, whereas the unmarked $\rho_k$ in the last
line is never declared a relation.

Process factor number from four down to one, using a fixed PBW order
and decreasing the cut index at each prescribed occurrence.  Every
commutator correction has fewer factors.  In a three-factor row with two ordinary
$F_1$-entries put the mark on the left, use the polarized transfer
\eqref{eq:polarize} below once, and append its lower-factor outputs.
In a $22$ row put the
ordinary factor on the left by $\rho Y=Y\rho+[\rho,Y]$, retaining that
placement thereafter.  Order a row by factor number, its unresolved
compulsory-transfer flag, its marked cut index (put $-1$ if unmarked),
and its number of PBW inversions, all in descending processing order.
Every same-factor output lowers the first datum at which it changes,
and every commutator output has fewer factors.  Thus the calculation
terminates.  At the
one-factor frontier, pieces with one label are first reassembled to
the full relation and only then set equal to zero in $K$.

At three factors replace the marked entry by
$\rho^{[2]}\in S_1$, project the other two entries to $S_0$, and take
the symmetric symbol.  The signed sum is
$\Xi_3\in K_3(S)_1$.  At two factors, after adjoining all corrections
from the preceding step, use $\rho^{[3]}\in R_1$ and project the
ordinary entry to $R_0$; call the result
$\Xi_2\in K_2(R)_1$.  The asserted kernel memberships hold because
$\rho$ maps to zero and the omitted pieces have weight at least three,
respectively four.  PBW identifies the factor-number symbol with a
symmetric algebra, so the factor-preserving part of
\eqref{eq:collect} is symmetric multiplication.  Hence, in weights at
most four, $\partial\Xi_3$ and $\partial\Xi_2$ are exactly the three-
and two-factor PBW coefficients of $\theta$.
Total bracket weight is preserved by \eqref{eq:collect}, and
\begin{equation}
 \aug(F)^5=\operatorname{span}_{\R}
 \{\text{PBW monomials of total weight at least }5\}.   \label{eq:PBWweight}
\end{equation}
Since $f$ has one factor,
\eqref{eq:witness}--\eqref{eq:PBWweight} imply
\begin{equation}
       \partial\Xi_3,\ \partial\Xi_2
       \quad\hbox{have weight at least five}.           \label{eq:cycles}
\end{equation}

For clarity, we compute the complete one-factor output of this
triangular collection.  Positive weights summing to at most four give
the exhaustive table
\begin{equation}
\begin{array}{c|c|c}
\text{factors}&\text{weight partitions}&\text{non-full exposed edge}\\ \hline
4&1111&\text{none}\\
3&111,112&211\\
2&11,12,13,22&31\text{ and the common }22\text{ root}\\
1&1,2,3,4&\text{primitive coefficient.}
\end{array}                                                   \label{eq:table}
\end{equation}
Indeed, a primitive descendant contains one homogeneous relation
component of weight $s$ and external entries of total weight $4-s$.
Replacing that component by its full labelled relation leaves only
the successive tail transitions
\[
       (1\longrightarrow2)+1+1,\qquad
       (1,2\longrightarrow3)+1,\qquad
       (1,2,3\longrightarrow4).
\]
These are $211$, $31$, and terminal reassembly.  If an external entry
has weight two, the sole additional possibility is $4=2+2$, the common
root treated below.  To see directly that no spectator term has been
lost, repeated transfer is the literal identity
\begin{equation}
 Bx_1\cdots x_t=
 \sum_{I\subseteq\{1,\ldots,t\}}x_{I^c}[B;x_I],         \label{eq:subsets}
\end{equation}
where the two orders are inherited from $1<\cdots<t$; induction on
$t$ proves the formula.  Applied simultaneously to all components of a
labelled relation, each proper-subset term, with its spectators, is a
component of the full labelled commutator $[\rho;x_I]$.  Only the full
subset can reach one factor.  This proves the asserted exhaustiveness.

For the $211$ transition write its relation as
$\rho=X-B+\rho_{\ge3}$.  Exact collection gives
\begin{equation}
 Bxy=xyB+x[B,y]+y[B,x]+[[B,x],y].                       \label{eq:Bxy}
\end{equation}
The primitive tooth of the tail $-B$ is $-[[B,x],y]$.
Its symmetric part is $-\Theta^+_\chi(B;xy)$, exactly
\eqref{eq:dT3}, because
$U_B=-B$.  Nothing is lost by symmetrizing: with $Y=[x,y]$,
\begin{equation}
 \rho xy=\tfrac12\rho(xy+yx)+\tfrac12Y\rho
                              +\tfrac12[\rho,Y].         \label{eq:polarize}
\end{equation}
The difference between the oriented and symmetric teeth is
$-\frac12[B,Y]$ by \eqref{eq:jacobi}.  The weight-four face of the
last output in \eqref{eq:polarize} is also $-\frac12[B,Y]$, with the
opposite boundary orientation,
because
$[\rho,Y]=[X,Y]-[B,Y]+[\rho_{\ge3},Y]$ is a full relation.
Thus these faces cancel.  The placed output $Y\rho$ is the same $22$
edge; the identity $\rho Y=Y\rho+[\rho,Y]$ shows that its other
orientation differs by that same full relation.  Hence the root is
counted once and contributes no additional primitive term.
This is labelwise for every $22$ row, not merely after summing their
symmetric symbols.  Since $F_2$ is spanned by brackets, reduce to
$Y=[x,y]$ and apply \eqref{eq:polarize} to
$\rho Y=\rho xy-\rho yx$.  The two symmetric three-factor terms cancel
and the remaining terms are exactly
$Y\rho+[\rho,Y]=\rho Y$.  Thus each occurrence carries its own closing
full-relation cell.

For the $31$ transition write the corresponding element of $R_1$ with
weight-three coordinate $-W$.  Moving this tail past $x$ gives
$-[W,x]=[x,W]$, whose character value is
$-\chi g(x,-W)=(dT_2)(Y;x)$ by \eqref{eq:dT2}.  The last transition
reassembles the full relation and is zero in $K$.  At the terminal
frontier all labelled terms therefore vanish; by
\eqref{eq:witness} the remaining one-factor coefficient is $f$, which
maps to $a$.  Consequently, with all occurrence signs retained,
the complete one-factor readout is the exact equality
\begin{equation}
             \chi(a)=dT_3(\Xi_3)+dT_2(\Xi_2).          \label{eq:readout}
\end{equation}
There is no confluence assumption here: use any fixed order in
\eqref{eq:collect},
and replace one occurrence at a time.  At every replacement, input
minus outputs is one of \eqref{eq:collect}, \eqref{eq:subsets}, or
\eqref{eq:polarize}; their sum
telescopes, and corrections strictly lower factor number.  Thus
the displayed transitions and \eqref{eq:readout} include every
labelled occurrence exactly once, while an
individual homogeneous tail has nowhere been declared a relation.
For comparison, all possible adjacent overlaps are also closed:
cuts commute by linearity and with transfers by the projection identity
following \eqref{eq:collect}; ordinary three-factor transfers give
Jacobi, disjoint transfers commute, and a marked overlap is
\[
 [u,[v,\rho]]-[v,[u,\rho]]=[[u,v],\rho].
\]
Interchanging $x,y$ across the polarized step gives precisely
$\rho(xy-yx)-[x,y]\rho-[\rho,[x,y]]=0$, the $22$ closure above.
Since every rule is an adjacent transfer, a cut, or this single
polarization, there are no further critical pairs.

Finally, \eqref{eq:dT3}--\eqref{eq:dT2}, \eqref{eq:cycles}, and the
weight-four support give
\[
 \chi(a)=T_3(\partial\Xi_3)+T_2(\partial\Xi_2)=0,
\]
contrary to the choice of $\chi$.  This proves the lemma.
\end{proof}

\section{Localization and conclusion}

\begin{corollary}\label{cor:class4}
If $K$ is an $\R$-Lie algebra with $\gamma_5(K)=0$, then
$\delta_5(K)=0$.
\end{corollary}

\begin{proof}
For $x\in\delta_5(K)$, naturality puts its image in
$\delta_5(K/K'')$.  It is killed by $2$ by \eqref{eq:PS}, since this
metabelian algebra has zero fifth lower-central term.
As $2$ is invertible, $x\in K''$; Lemma~\ref{lem:packet} gives $x=0$.
\end{proof}

\begin{theorem}
For every Lie ring $L$, the abelian group
\[
                         \delta_5(L)/\gamma_5(L)
\]
has no nonzero element of odd order.
\end{theorem}

\begin{proof}
Suppose that an odd-order class is represented by $x\in\delta_5(L)$,
and put $M=L/\gamma_5(L)$.  Its image $\bar x$ is nonzero and belongs
to $\delta_5(M)$.  Localize at $2$:
\[
                            K=\R\otimes_\Z M.
\]
The element $1\otimes\bar x$ is still nonzero: the kernel of
$M\to K$ consists of elements killed by a power of $2$.  Base change
for enveloping algebras gives
\[
 \U_\R(K)=\R\otimes_\Z\U_\Z(M),\qquad
 \aug_\R(K)^5=\R\otimes_\Z\aug_\Z(M)^5,
\]
so $1\otimes\bar x\in\delta_5(K)$.  But $K$ has class at most four,
and Corollary~\ref{cor:class4} says that $\delta_5(K)=0$, a
contradiction.
\end{proof}

\begin{remark}
Only the averages in \eqref{eq:theta} and \eqref{eq:polarize} require
a denominator, and that denominator is $2$.  Formula
\eqref{eq:polarize} is the small point of the proof: it
places the $22$ term and the full commutator relation in the same
identity, so the Jacobi root \eqref{eq:jacobi} cannot be omitted or
counted twice.
\end{remark}

\begin{thebibliography}{9}

\bibitem{PS}
I.~B.~S. Passi and T.~Sicking,
\emph{Dimension quotients of metabelian Lie rings},
Internat. J. Algebra Comput. \textbf{27} (2017), 251--258;
\href{https://arxiv.org/abs/1602.04919}{arXiv:1602.04919}.

\end{thebibliography}

\end{document}
